\inicializarSeccion{Simulación en MATLAB}

\inicializarSubseccion{Descripción de la simulación}
\tab La simulación que producimos en el desarrollo de esta etapa fue el gráfico que produce una linea de corriente eléctrica y una espira de corriente utilizando la ley de Biot-Savart; apoyandonos de resultados clásicos que salen de la ley de Ampere. Para esto, se decidio realizar la simulación en $\mathbb{R^3}$ para tener una mejor noción acerca de como se comporta el campo para la realización correcta de la simulación de la torre de caída libre. 

\inicializarSubseccion{Metodología}
\tab Antes de iniciar con la planeación del código, se realizo un diagrama de flujo para tener una mejor planeación de la simulación. Dicho diagrama de flujo se muestra en la \ref{fig:diagramaDeFlujo}, figura que se mostrará a continuación:

\begin{figure}[H]
    \centering
    \includegraphics[width=0.5\textwidth]{\nat{parte1/diagramaDeFlujo.png}}
    \caption{Diagrama de flujo de la simulación inicial.}
    \label{fig:diagramaDeFlujo}
\end{figure}

\tab Para la simulación se decidió utilizar la función \texttt{quiver3} de MATLAB, la cual nos ayuda a graficar campos vectoriales en $\mathbb{R}^3$, esto con el objetivo de poder visualizar el campo magnético que se genera por medio de una corriente eléctrica. Para esto, se utilizó la ley de Biot-Savart para calcular el campo magnético generado por una linea de corriente y por una espira circular.

\inicializarSubseccion{Herramientas técnicas}
\tab Como se pudo observar en el diagrama de flujo, se utilizó muchisimos conceptos físicos para la realización del simulador. Por mencionar algunos, tenemos la ley de Biot-Savart, la ley de Ampere, las corrientes de Eddy, la parametrización de la espira circular, entre otros. Sin embargo, para la realización del código en general, se necesitaron funciones especiales como \texttt{quiver3} para poder colocar las flechas que representan el campo magnético. También se utilizó la función \texttt{linspace} para generar los puntos en los cuales se va a evaluar el campo magnético. Se utilizó la función \texttt{meshgrid} para generar una malla de puntos en $\mathbb{R}^3$ para poder graficar el campo magnético en todo el espacio. Por último, se utilizó el método \texttt{reshape} para poder acotar las matrices que se generan por medio de la ley de Biot-Savart y poder graficar el campo magnético de una manera más clara. Utilizando estas herramientas técnicas, se logró generar una simulación que muestra el campo magnético generado por una linea de corriente y por una espira circular, lo cual nos ayuda a entender mejor el comportamiento del campo magnético en la torre de caída libre.


\inicializarSubseccion{Resultados iniciales}
\tab En esta sección se muestran los resultados obtenidos por medio de la simulación. La \autoref{fig:inicial} y \autoref{fig:vistaTransversal} muestran la simulación del campo generado por una linea de corriente y por una espira circular.

\begin{figure}[H]
    \begin{minipage}[t]{0.48\textwidth}
        \centering
        \includegraphics[width=1\textwidth]{\nat{parte1/inicial.png}}
        \caption{Vista inicial de la simulación.}
        \label{fig:inicial}
    \end{minipage}\hfill
    \begin{minipage}[t]{0.48\textwidth}
        \centering
        \includegraphics[width=1\textwidth]{\nat{parte1/vistaTransversal.png}}
        \caption{Vista sagital de la simulación.}
        \label{fig:vistaTransversal}
    \end{minipage}
\end{figure}

\tab En las \autoref{fig:vistaxy} y \autoref{fig:curvasDeNivel} se muestran la vista en el plano $xy$ de la simulación y las curvas de nivel del campo magnético respectivamente..

\begin{figure}[H]
    \begin{minipage}[t]{0.48\textwidth}
        \centering
        \includegraphics[width=1\textwidth]{\nat{parte1/vistaxy.png}}
        \caption{Vista en el plano XY de la simulación.}
        \label{fig:vistaxy}
    \end{minipage}\hfill
    \begin{minipage}[t]{0.48\textwidth}
        \centering
        \includegraphics[width=1\textwidth]{\nat{parte1/curvasDeNivel.png}}
        \caption{Curvas de nivel en el plano $xy$ para el campo magnético.}
        \label{fig:curvasDeNivel}
    \end{minipage}
\end{figure}

\inicializarSubseccion{Observaciones y comportamiento de la góndola con el frenado magnético}

Estos resultados son precisamente los esperados para la simulación; imitando el resultado experimental que se obtiene por medio de la ley de Ampere. Con esta base, tenemos las herramientas necesarias para poder realizar la simulación de la torre de caída libre con frenado magnético. En este caso, y utilizando este modelo, podemos atribuir otro diagrama de flujo que conecta este nodo de electromagnetico con un nodo de caida libre; combinandolos en un solo nodo magneto mecánico. Tal y como se muestra en la \autoref{fig:diagramaDeFlujoII}, figura que se muestra a continuación:

\begin{figure}[H]
    \centering
    \includegraphics[width=0.5\textwidth]{\nat{parte2/diagramaRetoII.png}}
    \caption{Diagrama de flujo de la simulación de la góndola.}
    \label{fig:diagramaDeFlujoII}
\end{figure}

Con este diagrama, podemos simular el comportamiento de la góndola en caída libre con frenado magnético, utilizando las herramientas y conceptos que se han desarrollado en esta etapa. Esto nos permitirá entender mejor el funcionamiento del sistema de frenado magnético y cómo se puede optimizar para mejorar su rendimiento.

\begin{figure}[H]
    \begin{minipage}[t]{0.48\textwidth}
        \centering
        \includegraphics[width=1\textwidth]{\nat{parte2/gondrolaInicial.png}}
        \caption{Vista de la góndola inicial.}
        \label{fig:gondrolaInicial}
    \end{minipage}\hfill
    \begin{minipage}[t]{0.48\textwidth}
        \centering
        \includegraphics[width=1\textwidth]{\nat{parte2/gondrolaEnDecaimiento.png}}
        \caption{Vista de la góndola cayendo.}
        \label{fig:gondrolaEnDecaimiento}
    \end{minipage}
\end{figure}

Con base a este modelo, podemos ver que la góndola experimenta una fuerza de frenado magnético que se opone a la fuerza de gravedad, lo que hace que la góndola caiga de manera más lenta. Para ver que tan lenta es, podemos ver las \autoref{fig:force_vs_position}, \autoref{fig:force_vs_time} y \autoref{fig:z_vs_time}, las cuales muestran la fuerza de frenado magnético en función de la posición, la velocidad en función del tiempo y la posición en función del tiempo respectivamente.
\begin{figure}[H]
    \begin{minipage}[t]{0.48\textwidth}
        \centering
        \includegraphics[width=1\textwidth]{\nat{parte2/force_vs_position.png}}
        \caption{Fuerza de frenado magnético en \\ función de la posición.}
        \label{fig:force_vs_position}
    \end{minipage}\hfill
    \begin{minipage}[t]{0.48\textwidth}
        \centering
        \includegraphics[width=1\textwidth]{\nat{parte2/force_vs_time.png}}
        \caption{Fuerza de frenado magnético \\ en función del tiempo.}
        \label{fig:force_vs_time}
    \end{minipage}
\end{figure}

\begin{figure}[H]
    \centering
    \includegraphics[width=0.5\textwidth]{\nat{parte2/z_vs_time.png}}
    \caption{Posición de la góndola en función del tiempo.}
    \label{fig:z_vs_time}
\end{figure}

Con estos resultados, podemos ver que la góndola cae de manera más lenta debido a la fuerza de frenado magnético, lo que demuestra el funcionamiento del sistema de frenado magnético. Además, podemos ver que la fuerza de frenado magnético aumenta a medida que la góndola se acerca al imán, lo que hace que la góndola caiga aún más lentamente. 

\tab Contestando algunas preguntas básicas acerca del comportamiento de la góndola. El juego mecanico tiene una velocidad terminal cuando este llega cayendo siete metros. Ya que a esta altura, empieza a afectar la fuerza del frenado magnetico; desacelerando la caída de la góndola. Para que esta velocidad terminal ocurra, solo tenemos que no acceder a la fuerza de frenado magnético, lo cual se puede lograr con un imán más pequeño. Lo que sucede y vimos en nuestra simulación fue que en el momento en el que la góndola cruza el anillo. El campo magnético cambia de dirección. Esto en casos en donde el dipolo magnético es muy fuerte puede llegar a que la góndola se quede atrapada en el anillo, lo cuál, evidentemente, es algo que queremos evitar. Con la fuerza predeterminada que se tiene, la góndola no se queda atrapada en el anillo, pero si se ve afectada por el cambio de dirección del campo magnético, lo que hace que la góndola caiga de manera más lenta.

\inicializarSubseccion{FEM inducida en la torre de caída libre}

\tab Sin embargo, esto no es todo lo que la góndola demuestra, y con este fundamento, queremos ahora pasar a la realización del cálculo de la FEM inducoda en toda la torre de caída libre, lo cual nos ayudará a entender mejor el comportamiento del sistema de frenado magnético y cómo se puede optimizar para mejorar su rendimiento. Para la realización del cálculo, se crearon dos nodos a nuestra instancia principal. El primer nodo se encarga de calcular numéricamente la FEM inducida a partir de la integración del flujo magnético antes, durante y después de que la góndola atraviese el anillo. El segundo nodo se encarga de graficar la FEM inducida en función de la posición, ademas de calcular la corriente inducida en un circuito RL y graficar esta corriente en función del tiempo. Con esto, se busca analizar la relación existente entre la posición de la góndola y la FEM inducida, así como también entender cómo es que esta FEM inducida afecta el movimiento de la góndola. Algo importante a mencionar, es que cambio el modelo del anillo a uno mas preciso; incluyendo multiples espiras (10 a una diferencia de altura de 0.25 m cada una) para poder generar un campo magnético más realista, lo cual nos ayuda a entender mejor el comportamiento del sistema de frenado magnético.

\tab Para la realización de esta sección, se optó por mover el anillo a 10 metros de altura por arriba de lo que se tenía antes, esto con objetivo de poder observar mejor el comportamiento de la FEM inducida cuando pasa la góndola por el anillo.

\begin{figure}
    \centering
    \includegraphics[width=0.45\textwidth]{\nat{parte3/diezMetros.png}}
    \caption{Góndola con los anillos iniciando ahora en $z=12$ m.}
    \label{fig:diezMetros}
\end{figure}

\tab Una manera de ver las consecuencias de este cambio es ver la gráfica de la fuerza de frenado magnético en función de la posición, lo cual se muestra en la \autoref{fig:force_vs_position} y la \autoref{fig:force_vs_time}, figuras que se muestra a continuación:

\begin{figure}[H]
    \begin{minipage}[t]{0.45\textwidth}
        \centering
        \includegraphics[width=1\textwidth]{\nat{parte3/force_vs_position.png}}
        \caption{Fuerza de frenado magnético en función de la posición con los anillos a 10 metros.}
        \label{fig:force_vs_position}
    \end{minipage}\hfill
    \begin{minipage}[t]{0.45\textwidth}
        \centering
        \includegraphics[width=1\textwidth]{\nat{parte3/force_vs_time.png}}
        \caption{Fuerza de frenado magnético en función del tiempo con los anillos a 10 metros.}
        \label{fig:force_vs_time}
    \end{minipage}
\end{figure}

\tab Con esto, se logró observar que la FEM inducida tiene un comportamiento oscilatorio, lo cual se debe a que el flujo magnético que atraviesa el anillo cambia de dirección a medida que la góndola se acerca al anillo, lo que hace que la FEM inducida también cambie de dirección. Tal y como lo demuestra la \autoref{fig:fem_vs_position}, la cual muestra la FEM inducida en función de la posición de la góndola.

\begin{figure}
    \centering
    \includegraphics[width=0.5\textwidth]{\nat{parte3/fem_vs_position.png}}
    \caption{FEM inducida en función de la posición de la góndola.}
    \label{fig:fem_vs_position}
\end{figure}

\tab Además, se logró observar que la corriente inducida en un circuito RL también tiene un comportamiento oscilatorio, lo cual se debe a que la FEM inducida también tiene un comportamiento oscilatorio. Tal y como lo demuestra la \autoref{fig:current_vs_position}, la cual muestra la corriente inducida en función de la posición de la góndola.

\begin{figure}
    \centering
    \includegraphics[width=0.5\textwidth]{\nat{parte3/current_vs_position.png}}
    \caption{Corriente inducida en función de la posición de la góndola.}
    \label{fig:current_vs_position}
\end{figure}

\tab Con estos resultados, se logró entender mejor el comportamiento del sistema de frenado magnético que ocurre en una torre de caída libre. Dicho analisis nos permitió observar que cuando el dipolo magnético es *debíl* y un objeto magnetico cae hacia el anillo. La FEM inducida tendrá un comportamiento oscilatorio simple. Esta observación que logramos realizarla dentro de nuestro simulador es un resultado clásico que se obtiene por medio de la ley de Faraday, lo cual nos ayuda a entender mejor el comportamiento del sistema de frenado magnético y cómo se puede optimizar para mejorar su rendimiento.